\documentclass[11pt,a4paper]{article}

%% ─── PACKAGES ───────────────────────────────────────────────────────────────
\usepackage[utf8]{inputenc}
\usepackage[T1]{fontenc}
\usepackage{amsmath, amssymb, amsfonts, amsthm}
\usepackage{mathtools}
\usepackage{geometry}
\usepackage{hyperref}
\usepackage{listings}
\usepackage{xcolor}
\usepackage{booktabs}
\usepackage{enumitem}
\usepackage{thmtools}
\usepackage{mdframed}
\usepackage{microtype}
\usepackage{cite}
\usepackage{bm}
\usepackage{bbm}
\usepackage{tikz}
\usetikzlibrary{arrows.meta, positioning, shapes.geometric}

\geometry{margin=1in, top=1.1in, bottom=1.1in}

%% ─── COLOURS ────────────────────────────────────────────────────────────────
\definecolor{codegreen}{rgb}{0,0.55,0.27}
\definecolor{codegray}{rgb}{0.45,0.45,0.45}
\definecolor{codepurple}{rgb}{0.50,0.0,0.70}
\definecolor{backcolour}{rgb}{0.96,0.96,0.94}
\definecolor{teal}{rgb}{0.0,0.41,0.44}
\definecolor{darkred}{rgb}{0.60,0.10,0.10}
\definecolor{boxbg}{rgb}{0.97,0.97,0.93}
\definecolor{boxrule}{rgb}{0.0,0.41,0.44}

%% ─── CODE STYLE ─────────────────────────────────────────────────────────────
\lstset{
  backgroundcolor=\color{backcolour},
  commentstyle=\color{codegreen}\itshape,
  keywordstyle=\color{codepurple}\bfseries,
  numberstyle=\tiny\color{codegray},
  stringstyle=\color{teal},
  basicstyle=\ttfamily\footnotesize,
  breaklines=true,
  captionpos=b,
  keepspaces=true,
  numbers=left,
  numbersep=6pt,
  showspaces=false,
  showstringspaces=false,
  tabsize=2,
  frame=single,
  rulecolor=\color{codegray!40},
  xleftmargin=1em,
  framexleftmargin=1em,
}

%% ─── THEOREM ENVIRONMENTS ───────────────────────────────────────────────────
\newtheoremstyle{tealstyle}
  {8pt}{6pt}{\normalfont}{0pt}{\bfseries\color{teal}}{.}{0.5em}{}

\theoremstyle{tealstyle}
\newtheorem{theorem}{Theorem}[section]
\newtheorem{lemma}[theorem]{Lemma}
\newtheorem{proposition}[theorem]{Proposition}
\newtheorem{corollary}[theorem]{Corollary}

\theoremstyle{definition}
\newtheorem{definition}[theorem]{Definition}
\newtheorem{remark}[theorem]{Remark}
\newtheorem{example}[theorem]{Example}

%% ─── BOXED ENVIRONMENT ──────────────────────────────────────────────────────
\newmdenv[
  backgroundcolor=boxbg,
  linecolor=boxrule,
  linewidth=1.2pt,
  leftmargin=0pt,
  rightmargin=0pt,
  innertopmargin=8pt,
  innerbottommargin=8pt,
  innerleftmargin=12pt,
  innerrightmargin=12pt,
  skipabove=10pt,
  skipbelow=6pt,
]{tealbox}

%% ─── OPERATORS & MACROS ─────────────────────────────────────────────────────
\newcommand{\Hilb}{\mathcal{H}}
\newcommand{\norm}[1]{\left\lVert #1 \right\rVert}
\newcommand{\abs}[1]{\left\lvert #1 \right\rvert}
\newcommand{\ip}[2]{\left\langle #1,\, #2 \right\rangle}
\newcommand{\R}{\mathbb{R}}
\newcommand{\N}{\mathbb{N}}
\newcommand{\Primes}{\mathbb{P}}
\newcommand{\eps}{\varepsilon}
\newcommand{\Lip}{\mathrm{Lip}}
\newcommand{\cond}{\kappa}
\newcommand{\spectrad}{\rho}
\newcommand{\topk}{C_k}
\newcommand{\qwit}{q}
\newcommand{\Tau}{\tau}
\newcommand{\Rmax}{R_{\max}}
\newcommand{\GovScore}{\mathtt{GovernanceScore}}
\newcommand{\HCalcIn}{\mathtt{HCalcInput}}
\newcommand{\TermState}{\mathcal{T}}
\newcommand{\F}{F}
\newcommand{\boolset}{\{0,1\}}
\newcommand{\converged}{\textsc{Converged}}
\newcommand{\freeze}{\textsc{Freeze}}
\newcommand{\freezecs}{\textsc{Freeze\textsubscript{CS}}}
\newcommand{\decoupled}{\textsc{Decoupled}}
\newcommand{\inconclusive}{\textsc{Inconclusive}}
\newcommand{\floor}[1]{\left\lfloor #1 \right\rfloor}
\newcommand{\sigmarank}{\sigma}
\newcommand{\dA}{\delta A}
\newcommand{\eigvec}{V}
\newcommand{\primed}{p}
\DeclareMathOperator{\diag}{diag}
\DeclareMathOperator{\spec}{spec}

%% ─── HYPERREF ────────────────────────────────────────────────────────────────
\hypersetup{
  colorlinks=true,
  linkcolor=teal,
  citecolor=teal,
  urlcolor=teal,
  pdftitle={HCALC Mathematical Appendix},
  pdfauthor={Citizen Gardens Engineering Team},
}

%% ─── TITLE ───────────────────────────────────────────────────────────────────
\title{%
  \textbf{Health-Calculator (HCALC)}\\[4pt]
  \large A Multi-Modality Convergence and Contraction Engine\\[6pt]
}
\author{%
  Citizen Gardens Engineering Team\\
  \small Phase Mirror · \texttt{packages/healthcare/hcalc}\\
  \small \href{https://github.com/MultiplicityFoundation/PhaseMirror-HQ}%
              {github.com/MultiplicityFoundation/PhaseMirror-HQ}
}
\date{April 2026\quad\textbf{[Prior Art Defensive Publication]}}

%% ═══════════════════════════════════════════════════════════════════════════
\begin{document}
\maketitle
\tableofcontents
\bigskip

\begin{tealbox}
\textbf{Scope of this Appendix.}\;
This document presents the complete mathematical derivations underlying the
H-Calculator (\texttt{hcalc}) framework.
Every theorem, lemma, and bound stated here corresponds directly to
code artefacts in \texttt{packages/healthcare/hcalc/} and constitutes
explicit prior art.
Nothing in this appendix is speculative: each result is either already
implemented or is the formal justification for a design decision visible in
the source.
\end{tealbox}

%% ═══════════════════════════════════════════════════════════════════════════
\section{Hilbert-Space Setting}
\label{sec:setting}

\begin{definition}[Patient State Space]
\label{def:state-space}
Let $\Hilb = L^2(\Omega,\mu)$ be a separable, real Hilbert space of
\emph{patient state vectors}, where $\Omega$ is the physiological measurement
domain and $\mu$ is a reference probability measure.
Inner product and induced norm are denoted $\ip{\cdot}{\cdot}$ and
$\norm{\cdot}$, respectively.
\end{definition}

\begin{definition}[Modality Tensor and Direct-Sum Assembly]
\label{def:assembly}
Let $n \in \N$ with $n \le \Rmax$ (see Definition~\ref{def:rmax}).
For each modality index $i \in \{1,\dots,n\}$ assign a prime
$\primed_i \in \Primes$ according to the canonical injection
\[
  \primed\colon i \mapsto \text{the } i\text{-th prime},
  \quad \primed_1 = 2,\; \primed_2 = 3,\; \primed_3 = 5, \ldots
\]
Each modality contributes a vector $\bm{m}_i \in \Hilb$,
normalised by \texttt{build\_modality\_tensor} to a tuple of floats.
The \emph{joint state vector} is the direct sum
\[
  \bm{x} = \bigoplus_{i=1}^{n} \bm{m}_i \;\in\; \bigoplus_{i=1}^{n}\Hilb,
\]
with induced norm $\norm{\bm{x}}^2 = \sum_{i=1}^{n}\norm{\bm{m}_i}^2$.
The \emph{modality mass} is $m \coloneqq \norm{\bm{x}}$.
\end{definition}

\begin{definition}[Diagnostic Operator]
\label{def:operator}
The \emph{diagnostic operator}
$T\colon \bigoplus_{i=1}^n \Hilb \to \bigoplus_{i=1}^n \Hilb$
maps the joint state vector to an updated posterior state.
$T$ is assumed bounded and linear for the purposes of the spectral analysis;
the nonlinear Lipschitz extension is treated in
Section~\ref{sec:lipschitz}.
\end{definition}

%% ─────────────────────────────────────────────────────────────────────────────
\section{Contraction Witness: Definition and Fundamental Properties}
\label{sec:witness}

\subsection{Scalar Witness Formula}
\label{sec:witness-formula}

\begin{definition}[Contraction Witness]
\label{def:witness}
Given scalar parameters $\Tau > 0$ (step-size / mass-transfer coefficient),
$m = \norm{\bm{x}} \ge 0$ (modality mass), and
$L > 0$ (Lipschitz bound of $T$; cf.\ Definition~\ref{def:lipschitz}),
the \emph{contraction witness} is
\begin{equation}
  \label{eq:witness}
  \qwit \;=\; \frac{\Tau \cdot m}{L}.
\end{equation}
The implementation floors $L$ at $10^{-12}$ to prevent division by zero:
\[
  L_{\text{safe}} = \max(L,\;10^{-12}).
\]
This floor introduces a negligible perturbation
$\abs{\qwit - \qwit_{\text{safe}}} = O(10^{-12})$ for any clinically
realistic $\Tau$ and $m$.
\end{definition}

\begin{definition}[Contractivity Predicate]
\label{def:contractive}
Let $\eps \in (0,1)$ be the profile-specific convergence epsilon
(default $\eps = 0.01$, cf.\ \texttt{config.CONVERGENCE\_EPSILON}).
The joint assembly is \emph{contractive} iff
\begin{equation}
  \label{eq:contractive}
  \qwit < 1 - \eps.
\end{equation}
\end{definition}

\subsection{Banach Fixed-Point Interpretation}

\begin{theorem}[Convergence of the Diagnostic Operator]
\label{thm:banach}
Suppose $T$ is a self-map of a non-empty closed subset
$\mathcal{C} \subseteq \bigoplus_{i=1}^n \Hilb$ and that the contraction
witness $\qwit$ satisfies \eqref{eq:contractive}.
Then $T$ is a Banach contraction on $\mathcal{C}$, and for every
$\bm{x}_0 \in \mathcal{C}$ the Picard iterates
$\bm{x}_{t+1} = T(\bm{x}_t)$ converge to a unique fixed point
$\bm{x}^* \in \mathcal{C}$ with error bound
\begin{equation}
  \label{eq:picard-bound}
  \norm{\bm{x}_t - \bm{x}^*}
  \;\le\; \frac{\qwit^t}{1 - \qwit}\,\norm{\bm{x}_1 - \bm{x}_0}.
\end{equation}
\end{theorem}

\begin{proof}
Standard Banach Fixed-Point Theorem.
The contraction constant of $T$ is upper-bounded by $\qwit$ because,
for any $\bm{x}, \bm{y} \in \mathcal{C}$,
\[
  \norm{T(\bm{x}) - T(\bm{y})}
  \;\le\; L\,\norm{\bm{x} - \bm{y}}
  \;\le\; \qwit\,\norm{\bm{x} - \bm{y}},
\]
under the witness-calibrated contractive regime.
Since $\qwit < 1$ by \eqref{eq:contractive}, the claim follows.
The error bound \eqref{eq:picard-bound} is the standard a-priori estimate.
\end{proof}

\begin{corollary}[Clinical Interpretation]
\label{cor:clinical}
The terminal state $\converged$ is \emph{issued if and only if}
$\bm{x}_t$ has entered the basin of attraction of the unique fixed point
$\bm{x}^*$ and the top-$k$ concentration $\topk \ge \theta_c$ verifies that
probability mass is concentrated on a clinically actionable differential
(see Section~\ref{sec:topk}).
\end{corollary}

\subsection{Temporal Chaining and Drift Detection}
\label{sec:temporal}

\begin{definition}[Temporal Contraction Sequence]
\label{def:temporal}
Let $\{\bm{x}_t\}_{t \ge 0}$ be a sequence of patient state observations.
For each $t$, compute the witness $\qwit_t$ via \eqref{eq:witness}.
The \emph{q-delta} is
\[
  \Delta_t \;\coloneqq\; \qwit_t - \qwit_{t-1},
  \qquad t \ge 1,
\]
implemented as the field \texttt{q\_delta} in \texttt{ContractionWitness}.
\end{definition}

\begin{proposition}[Drift Detection via q-Delta]
\label{prop:drift}
Let $d_{\max} > 0$ be the profile parameter \texttt{drift\_bound\_max}.
If $\sum_{s=1}^{t} \Delta_s > d_{\max}$, the sequence $\{\qwit_t\}$ is
drifting toward instability and the system should emit \freeze\ or trigger
an alert before $\qwit_t \ge 1 - \eps$.
\end{proposition}

\begin{proof}
By telescoping, $\qwit_t = \qwit_0 + \sum_{s=1}^{t} \Delta_s$.
If this sum exceeds $d_{\max}$ while $\qwit_0 + d_{\max} < 1 - \eps$,
then the witness is still contractive but trending unstable.
The drift gate catches this before contractivity is lost.
\end{proof}

%% ─────────────────────────────────────────────────────────────────────────────
\section{Lipschitz Operator Norm Bounds}
\label{sec:lipschitz}

\begin{definition}[Lipschitz Bound]
\label{def:lipschitz}
The diagnostic operator $T$ is said to have \emph{Lipschitz constant} $L > 0$
if
\begin{equation}
  \label{eq:lipschitz}
  \norm{T(\bm{x}) - T(\bm{y})} \;\le\; L\,\norm{\bm{x} - \bm{y}}
  \quad \forall\; \bm{x}, \bm{y} \in \Hilb.
\end{equation}
For a bounded linear operator, $L = \norm{T}_{\mathrm{op}}$.
\end{definition}

\begin{lemma}[Operator Norm as Lipschitz Supremum]
\label{lem:op-norm}
For a bounded linear $T\colon \Hilb \to \Hilb$,
\[
  L \;=\; \norm{T}_{\mathrm{op}}
    \;=\; \sup_{\bm{x} \ne 0} \frac{\norm{T\bm{x}}}{\norm{\bm{x}}}
    \;=\; \sup_{\norm{\bm{x}}=1} \norm{T\bm{x}}.
\]
\end{lemma}

\begin{proof}
Direct from the definition of the operator norm.
\end{proof}

\begin{theorem}[Stability from Lipschitz Bound]
\label{thm:stability}
Let $\bm{x}$ be the current joint state vector and let
$T$ be the diagnostic operator with Lipschitz constant $L$.
Define the \emph{normalized witness ratio}
\[
  \rho^* \;\coloneqq\; \frac{\Tau \cdot m}{L}.
\]
Then:
\begin{enumerate}[label=(\roman*)]
  \item $\rho^* < 1 - \eps$ implies $T$ is a Banach contraction
    (Theorem~\ref{thm:banach}).
  \item $L < 1$ is a necessary condition for the Gate~1 stability check
    to pass.
  \item $L \ge 1$ alone, even with $\rho^* < 1 - \eps$, triggers \freeze\
    because the operator itself is expansive in the unit ball.
\end{enumerate}
\end{theorem}

\begin{proof}
\textit{(i)} Follows from Theorem~\ref{thm:banach}.
\textit{(ii,iii)} Even if $m$ is small enough that
$\rho^* = \Tau m / L < 1 - \eps$ when $L \ge 1$,
the operator is still non-contractive on the unit ball:
for $\norm{\bm{x}} = 1$, $\norm{T\bm{x}} \le L\norm{\bm{x}} = L \ge 1$.
Thus small mass is not a structural guarantee.
The separate Gate~1 check closes this gap.
\end{proof}

\begin{remark}
The condition $L < 1$ is equivalent to the spectral radius
$\spectrad(T) < 1$ when $T$ is normal.
For non-normal $T$, $\spectrad(T) \le \norm{T}_{\mathrm{op}} = L$,
so $L < 1$ is the stronger safe-side condition.
\end{remark}

%% ─────────────────────────────────────────────────────────────────────────────
\section{Spectral Risk and Bauer-Fike Bounds}
\label{sec:spectral}

\subsection{Eigenvalue Perturbation}

\begin{theorem}[Bauer-Fike Bound]
\label{thm:bauerfike}
Let $A \in \R^{n \times n}$ be diagonalisable as $A = \eigvec\,\Lambda\,\eigvec^{-1}$
where $\Lambda = \diag(\lambda_1,\dots,\lambda_n)$ and $\eigvec$ is the
matrix of right eigenvectors.
Let $\dA \in \R^{n \times n}$ be a perturbation and $\mu$ an eigenvalue of
$A + \dA$.
Then, for any $p$-norm,
\begin{equation}
  \label{eq:bauerfike}
  \min_{i} \abs{\lambda_i - \mu}
  \;\le\; \cond_p(\eigvec)\,\norm{\dA}_p.
\end{equation}
\end{theorem}

\begin{proof}
Standard Bauer-Fike.
Suppose $\mu$ is not an eigenvalue of $A$.
Then $(A - \mu I)$ is invertible, and
$(A + \dA - \mu I) = (A - \mu I)(I + (A - \mu I)^{-1}\dA)$.
For $\mu$ to be an eigenvalue of $A + \dA$ we need
$\norm{(A-\mu I)^{-1}\dA}_p \ge 1$.
Writing $(A - \mu I)^{-1} = \eigvec(\Lambda - \mu I)^{-1}\eigvec^{-1}$
and applying sub-multiplicativity yields
\[
  1 \;\le\; \cond_p(\eigvec)\cdot
        \frac{\norm{\dA}_p}{\min_i\abs{\lambda_i - \mu}},
\]
which rearranges to \eqref{eq:bauerfike}.
\end{proof}

\subsection{Runtime Cap from Condition Number Growth}

\begin{proposition}[Exponential Growth of Condition Number]
\label{prop:condgrowth}
For a block-diagonal operator
$A = \diag(A_1, \dots, A_n)$ assembled from $n$ modality blocks,
\[
  \cond_2(\eigvec) \;\le\; \prod_{i=1}^n \cond_2(\eigvec_i),
\]
where $\eigvec_i$ is the eigenvector matrix of block $A_i$.
In the worst case, $\cond_2(\eigvec_i) \ge c > 1$ for each $i$, giving
$\cond_2(\eigvec) = O(c^n)$.
\end{proposition}

\begin{proof}
The eigenvector matrix of a block-diagonal matrix is block-diagonal:
$\eigvec = \diag(\eigvec_1,\dots,\eigvec_n)$.
Hence
\[
\cond_2(\eigvec)
\le \prod_i \cond_2(\eigvec_i).
\]
The worst-case bound follows immediately.
\end{proof}

\begin{corollary}[Justification of $\Rmax = 16$]
\label{cor:rmax}
Let $\delta_{\max} > 0$ be the maximum clinically tolerable perturbation
of the eigenvalue spectrum.
From \eqref{eq:bauerfike}, we require
\[
  \cond_2(\eigvec)\,\norm{\dA}_2 \;\le\; \delta_{\max}.
\]
Substituting the exponential bound from Proposition~\ref{prop:condgrowth}
with a conservative $c = 1.5$ and $\norm{\dA}_2 = 0.01$,
the bound is exceeded at
\[
  n^* = \floor*{\frac{\log(\delta_{\max} / \norm{\dA}_2)}{\log c}}
      \approx 16.
\]
This justifies \texttt{R\_MAX\_RUNTIME = 16}.
\end{corollary}

\begin{definition}[$\Rmax$ Runtime Cap]
\label{def:rmax}
The \emph{runtime cap} $\Rmax \in \N$ is the maximum number of modalities
$n$ for which the Bauer-Fike spectral bound can be guaranteed within
clinical safety tolerances.
Gate~4 enforces $n \le \Rmax$; violation emits
$\decoupled$.
\end{definition}

\subsection{Spectral Radius Guard}

\begin{proposition}[Spectral Radius Bound]
\label{prop:spectrad}
Let $\lambda_{\max} = \max_i \abs{\lambda_i(T)}$ be the spectral radius of
the assembly operator.
The resonance gate enforces
\begin{equation}
  \label{eq:resonance}
  r_{\min} \;\le\; \spectrad(T) \;\le\; r_{\mathrm{safe}},
\end{equation}
where $r_{\min} = 0.10$ and $r_{\mathrm{safe}} = 0.90$.
Values below $r_{\min}$ indicate a degenerate operator;
values above $r_{\mathrm{safe}}$ approach the instability boundary.
\end{proposition}

%% ─────────────────────────────────────────────────────────────────────────────
\section{Prime-Indexed Modality Decomposition}
\label{sec:prime}

\subsection{Construction}

\begin{definition}[Prime-Indexed Eigenspace]
\label{def:prime-eigenspace}
Let $\primed_i$ be the prime assigned to modality $i$.
The \emph{prime-indexed eigenspace} $\mathcal{E}_{p_i}$ is the invariant
subspace of $T$ corresponding to the spectral component associated with
modality $i$.
The joint eigenvector matrix decomposes as
\[
  \eigvec \;=\; \bigoplus_{i=1}^n \eigvec_{\primed_i},
\]
where $\eigvec_{\primed_i}$ governs the $i$-th modality block exclusively.
\end{definition}

\begin{theorem}[Prime Orthogonality and Quarantine]
\label{thm:prime-quarantine}
Under Definition~\ref{def:prime-eigenspace}, the condition number of the
joint assembly satisfies
\begin{equation}
  \label{eq:prime-bound}
  \norm{\dA}_{\mathrm{joint}}
  \;\le\; \prod_{i=1}^n \cond_{\primed_i}(\eigvec_{\primed_i})^{-1},
\end{equation}
where $\cond_{\primed_i}$ denotes the condition number computed in the
$\ell^{\primed_i}$-norm.
Consequently, a single pathological modality $j$ with
$\cond_{\primed_j}(\eigvec_{\primed_j}) \gg 1$
decreases only the $j$-th factor in the product, leaving all other modality
eigenspaces unaffected.
\end{theorem}

\begin{proof}
Because $\eigvec = \bigoplus_i \eigvec_{\primed_i}$ is block-diagonal,
cross-block interference terms vanish.
The perturbation $\dA$ acts within each eigenspace independently, so
applying the Bauer-Fike estimate blockwise yields the stated product bound.
\end{proof}

\begin{corollary}[Fault Isolation]
\label{cor:fault-isolation}
If modality $j$ becomes unstable, Gate~4 may decouple modality $j$ from the
joint assembly.
All remaining modalities $i \ne j$ retain their convergence guarantees,
and the system may emit a \decoupled\ state with per-modality fallback
rather than a global \freeze.
\end{corollary}

%% ─────────────────────────────────────────────────────────────────────────────
\section{Top-$k$ Concentration}
\label{sec:topk}

\begin{definition}[GovernanceScore and Ranked Differential]
\label{def:gov-score}
A \emph{GovernanceScore} is a pair $(c, p) \in \Sigma \times [0,1]$,
where $\Sigma$ is a domain-specific code space and $p$ is the posterior
probability weight.
The \emph{ranked differential} is a finite tuple
$\bm{s} = (s_1,\dots,s_N) \in \GovScore^N$
sorted by probability in descending order.
\end{definition}

\begin{definition}[Top-$k$ Concentration]
\label{def:topk}
For $k \ge 1$ and ranked differential $\bm{s}$, the
\emph{top-$k$ concentration} is
\begin{equation}
  \label{eq:topk}
  \topk(\bm{s}) \;\coloneqq\; \sum_{i=1}^{\min(k,N)} p_{\sigmarank(i)}.
\end{equation}
\end{definition}

\begin{theorem}[Concentration Necessary for Convergence]
\label{thm:topk}
Let $\theta_c \in (0,1)$ be the profile concentration threshold
(default $\theta_c = 0.60$).
Even if $\qwit < 1 - \eps$, the terminal state is
$\converged$ \emph{only if} $\topk(\bm{s}) \ge \theta_c$.
\end{theorem}

\begin{proof}
This is an axiomatic architectural condition.
If $\topk < \theta_c$, the posterior mass remains too dispersed for
a safe action route.
Hence HCALC requires both dynamical convergence and epistemic concentration.
\end{proof}

\begin{lemma}[Monotonicity of $\topk$]
\label{lem:topk-mono}
$\topk(\bm{s})$ is non-decreasing in $k$ and satisfies
$\topk(\bm{s}) \le \topk[k+1](\bm{s}) \le 1$ for all $k < N$.
\end{lemma}

\begin{proof}
Immediate from \eqref{eq:topk}.
\end{proof}

%% ─────────────────────────────────────────────────────────────────────────────
\section{Four-Gate Decision Function}
\label{sec:gates}

\subsection{Gate Definitions}

\begin{definition}[Gate Predicates]
\label{def:gates}
Define gate predicates $g_i \in \boolset$ as follows:
\[
\begin{aligned}
  g_1 &= \mathbbm{1}\bigl[\texttt{physio\_bounds\_ok}
         \wedge L < 1
         \wedge \texttt{drift} \le d_{\max}
         \wedge r_{\min} \le \texttt{resonance} \le r_{\mathrm{safe}}\bigr], \\
  g_2 &= \mathbbm{1}[\qwit < 1 - \eps], \\
  g_3 &= \mathbbm{1}[\texttt{epistemic\_verified}], \\
  g_4 &= \mathbbm{1}[n \le \Rmax].
\end{aligned}
\]
\end{definition}

\subsection{Formal Decision Function}

\begin{definition}[Terminal State Decision Function]
\label{def:decision}
The \emph{terminal state decision function}
$\F\colon \boolset^4 \to \TermState$ is defined by the ordered
case analysis:
\begin{equation}
  \label{eq:decision}
  \F(g_1, g_2, g_3, g_4) \;=\;
  \begin{cases}
    \decoupled  & \text{if } \neg g_4, \\[4pt]
    \freezecs   & \text{if } g_4 \wedge g_1 \wedge \neg g_2, \\[4pt]
    \freeze     & \text{if } g_4 \wedge \neg g_1, \\[4pt]
    \inconclusive
                & \text{if } g_4 \wedge g_1 \wedge g_2 \wedge \neg g_3, \\[4pt]
    \inconclusive
                & \text{if } g_4 \wedge g_1 \wedge g_2 \wedge g_3
                             \wedge \topk < \theta_c, \\[4pt]
    \converged  & \text{if } g_4 \wedge g_1 \wedge g_2 \wedge g_3
                             \wedge \topk \ge \theta_c.
  \end{cases}
\end{equation}
\end{definition}

\begin{theorem}[Totality and Determinism of $\F$]
\label{thm:total}
$\F$ is total and deterministic.
No input admits more than one output.
\end{theorem}

\begin{proof}
There are $2^4 = 16$ possible input combinations.
The case analysis in \eqref{eq:decision} is exhaustive and priority-ordered,
so every Boolean input has exactly one terminal state.
\end{proof}

\begin{theorem}[Non-Commutativity of Gate Ordering]
\label{thm:noncommute}
The mapping $(g_1,g_2,g_3,g_4) \mapsto \F(g_1,g_2,g_3,g_4)$ is
not symmetric in its arguments: swapping the evaluation order
changes the output for at least one input combination.
\end{theorem}

\begin{proof}
Consider $(g_1,g_2,g_3,g_4) = (1,0,0,1)$.
Under \eqref{eq:decision}, this yields $\freezecs$.
If Gate~3 were evaluated first, the same input could yield
$\inconclusive$.
Hence ordering is structurally non-commutative.
\end{proof}

\begin{remark}
This formalises the architectural invariant that cytokine storm has higher
priority than mere provenance failure.
\end{remark}

%% ─────────────────────────────────────────────────────────────────────────────
\section{Cytokine Storm as a Universal Instability Archetype}
\label{sec:cytokine}

\begin{definition}[Cytokine Storm Failure Topology]
\label{def:cytokine}
A system exhibits the \emph{cytokine storm failure topology} when:
\begin{enumerate}[label=(\roman*)]
  \item Local indicators satisfy nominal bounds ($g_1 = 1$), yet
  \item The joint operator is non-contractive ($g_2 = 0$).
\end{enumerate}
\end{definition}

\begin{proposition}[Cross-Domain Instances]
\label{prop:crossdomain}
The cytokine storm topology generalises across domains:
\[
\begin{array}{lll}
\toprule
\textbf{Domain} & \textbf{Gate 1 passes} & \textbf{Gate 2 fails} \\
\midrule
\text{Clinical (sepsis)} & \text{Vitals normal} & \text{Cytokine cascade diverging} \\
\text{Finance (bubble)} & \text{Asset prices rising} & \text{Leverage operator non-contractive} \\
\text{AI (hallucination)} & \text{Fluent output} & \text{Internal state diverging} \\
\text{Education (burnout)} & \text{Completion metrics high} & \text{Knowledge consolidation failing} \\
\bottomrule
\end{array}
\]
\end{proposition}

%% ─────────────────────────────────────────────────────────────────────────────
\section{Domain Profile: Parameterised Compliance Surface}
\label{sec:profiles}

\begin{definition}[DomainProfile as Compliance Surface]
\label{def:profile}
A \emph{DomainProfile} $\Pi$ is a frozen tuple
\[
  \Pi = (\Rmax,\; d_{\max},\; r_{\min},\; r_{\mathrm{safe}},\;
         k,\; \theta_c,\; \eps,\; c_{\min},\; \mathcal{A},\; \Phi),
\]
where $\mathcal{A}$ is the audit-trail flag and
$\Phi$ is a finite set of protected fields.
\end{definition}

\begin{theorem}[Profile Validity Invariant]
\label{thm:profile-valid}
A profile is valid iff all of the following hold:
\[
\begin{aligned}
  &\Rmax > 0, \quad d_{\max} > 0, \quad \eps > 0, \\
  &0 \le c_{\min} \le 1, \quad r_{\min} \ge 0, \quad r_{\mathrm{safe}} > r_{\min}, \\
  &k > 0, \quad 0 \le \theta_c \le 1.
\end{aligned}
\]
\end{theorem}

\begin{proof}
Necessity follows because each threshold must define a non-vacuous predicate.
Sufficiency follows because all gates are then well-defined and the terminal
state map remains total.
\end{proof}

%% ─────────────────────────────────────────────────────────────────────────────
\section{V(D)J Three-Key Activation: Cryptographic Governance}
\label{sec:vdj}

\begin{definition}[Three-Key Signature]
\label{def:threekey}
A \emph{three-key signature set} for proposal hash $h$ is
a set of triples
$(\mathtt{id}_j, \mathtt{role}_j, H_j)$ satisfying:
\begin{enumerate}[label=(\roman*)]
  \item Hash consistency.
  \item Role coverage over the required signatory roles.
  \item At least three unique signatories.
\end{enumerate}
\end{definition}

\begin{theorem}[Single-Actor Bypass Impossibility]
\label{thm:bypass}
Under Definition~\ref{def:threekey}, no single principal can produce
a valid three-key signature set alone.
\end{theorem}

\begin{proof}
Condition (iii) requires at least three distinct signatory identities.
A single principal contributes only one identity.
Therefore the condition fails.
\end{proof}

%% ─────────────────────────────────────────────────────────────────────────────
\section{OLHES Integration: Convergence as a Control-Theoretic Loop}
\label{sec:olhes}

\begin{definition}[OLHES Closed-Loop System]
\label{def:olhes}
The \emph{One-Loop Health-Education System} (OLHES) is the discrete
dynamical system
\[
  \bm{x}_{t+1} \;=\; G\bigl(A(\bm{x}_t),\; \F(\bm{g}_t)\bigr),
\]
where $A$ is the action operator,
$G$ is the outcome-feedback operator, and
$\bm{g}_t = (g_1^t, g_2^t, g_3^t, g_4^t)$ are the HCALC gate values.
\end{definition}

\begin{proposition}[HCALC as Homeostatic Checkpoint]
\label{prop:homeostatic}
HCALC implements a homeostatic checkpoint within OLHES:
action is applied only when $\F = \converged$,
which guarantees the system is in the basin of attraction of a unique
fixed point.
\end{proposition}

\begin{theorem}[Lyapunov Stability of the OLHES Loop Under HCALC Gating]
\label{thm:lyapunov}
Define the Lyapunov function $V_t = \norm{\bm{x}_t - \bm{x}^*}^2$.
Under HCALC gating, $V_{t+1} \le \qwit^2 V_t$ for all $t$ at which an
action is applied.
Consequently the loop is exponentially stable in expectation.
\end{theorem}

\begin{proof}
When $\F = \converged$, Theorem~\ref{thm:banach} gives
\[
\norm{\bm{x}_{t+1} - \bm{x}^*}
 = \norm{T(\bm{x}_t) - \bm{x}^*}
 \le \qwit \norm{\bm{x}_t - \bm{x}^*}.
\]
Squaring yields $V_{t+1} \le \qwit^2 V_t$.
Since $\qwit < 1$, the decay is exponential.
\end{proof}

%% ─────────────────────────────────────────────────────────────────────────────
\appendix
\section{Key Code Artefacts}
\label{app:code}

\subsection{\texttt{contraction.py} — Witness Computation}
\begin{lstlisting}[language=Python, caption={Contraction witness with temporal chaining.}]
@dataclass(frozen=True)
class ContractionWitness:
    q: float
    is_contractive: bool
    policy_code: str
    q_delta: float | None = None

def compute_witness(
    tau: float,
    mass: float,
    lipschitz_term: float,
    epsilon: float,
    prior_witness_q: float | None = None,
) -> ContractionWitness:
    denom = max(float(lipschitz_term), 1e-12)
    q = float(tau) * float(mass) / denom
    is_contractive = q < (1.0 - float(epsilon))
    q_delta = (q - prior_witness_q) if prior_witness_q is not None else None
    return ContractionWitness(
        q=q,
        is_contractive=is_contractive,
        policy_code="joint_contractive" if is_contractive else "joint_non_contractive",
        q_delta=q_delta,
    )
\end{lstlisting}

\subsection{\texttt{convergence.py} — Four-Gate Routing}
\begin{lstlisting}[language=Python, caption={Ordered non-commutative gate routing.}]
def route_terminal_state(
    *,
    physio_bounds_ok: bool,
    joint_contractive: bool,
    epistemic_verified: bool,
    runtime_cap_ok: bool,
    confidence_ok: bool,
    stability_ok: bool = True,
) -> TerminalState:
    if not runtime_cap_ok:
        return TerminalState.DECOUPLED
    if physio_bounds_ok and not joint_contractive:
        return TerminalState.FREEZE_CYTOKINE_STORM
    if not stability_ok:
        return TerminalState.FREEZE
    if (
        physio_bounds_ok
        and joint_contractive
        and epistemic_verified
        and runtime_cap_ok
        and confidence_ok
    ):
        return TerminalState.CONVERGED
    return TerminalState.INCONCLUSIVE
\end{lstlisting}

\subsection{\texttt{ranking.py} — Top-k Concentration}
\begin{lstlisting}[language=Python, caption={Top-k concentration helper.}]
def top_k_mass(scores: tuple[DiagnosisScore, ...], k: int = 3) -> float:
    ordered = sorted((item.probability for item in scores), reverse=True)
    return sum(ordered[: max(1, int(k))]) if ordered else 0.0
\end{lstlisting}

\appendix
\section*{Mathematical Appendix}

\section{Normed Structure on Histories and Operators}

In this appendix we make explicit the functional-analytic structure underlying the HCALC operators and state operator norm bounds and basic continuity results. [file:1]

\subsection{Spaces of Histories and Functionals}

Fix a finite time horizon \(T \in \mathbb{N}\). [file:1] Let \(S, A, O\) be finite sets of states, actions, and observations, respectively. [file:1] A history of length \(t \le T\) is a sequence
\[
  h_t = (s_0, a_0, o_0, s_1, a_1, o_1, \dots, s_t) \in H_t \subset (S \times A \times O)^{t} \times S,
\]
and we define the full history space
\[
  H := \bigcup_{t=0}^T H_t.
\] [file:1]

Let \(\mathcal{F} := \{ f : H \to \mathbb{R} \}\) be the set of real-valued functionals on histories. [file:1] We endow \(\mathcal{F}\) with the sup-norm
\[
  \|f\|_\infty := \sup_{h \in H} |f(h)|.
\] [file:1]

\begin{lemma}[Banach structure]
With the sup-norm, \((\mathcal{F}, \|\cdot\|_\infty)\) is a Banach space. [file:1]
\end{lemma}

\begin{proof}
The space of bounded real-valued functions on a set with the sup-norm is complete; this is standard. Since \(H\) is finite under the imposed horizon \(T\), every \(f \in \mathcal{F}\) is automatically bounded and completeness follows from the finite-dimensionality. [file:1]
\end{proof}

\subsection{Transition Kernel and History Extension Operator}

Let \(B\) be a Markov kernel giving the probability of next state and observation:
\[
  B : H \times A \to \Delta(S \times O), \qquad
  (h,a) \mapsto B(\cdot \mid h,a).
\] [file:1]

For fixed policy \(\pi : H \to A\), define the one-step \emph{history extension operator}
\[
  T_\pi : \mathcal{F} \to \mathcal{F}
\]
acting on a functional \(f \in \mathcal{F}\) by
\[
  (T_\pi f)(h) 
  := \mathbb{E}_{(s',o') \sim B(\cdot \mid h, \pi(h))}\bigl[f(h \cdot ( \pi(h), o', s'))\bigr],
\]
where \(h \cdot (\pi(h), o', s')\) denotes concatenation of one action–observation–state triple. [file:1]

\begin{lemma}[Linearity]
For any fixed policy \(\pi\), the operator \(T_\pi\) is linear on \(\mathcal{F}\). [file:1]
\end{lemma}

\begin{proof}
Let \(f,g \in \mathcal{F}\) and \(\alpha,\beta \in \mathbb{R}\). For all \(h \in H\),
\begin{align*}
  T_\pi(\alpha f + \beta g)(h)
  &= \mathbb{E}_{(s',o') \sim B(\cdot \mid h, \pi(h))}[(\alpha f + \beta g)(h \cdot (\pi(h),o',s'))] \\
  &= \alpha \mathbb{E}[f(h \cdot (\pi(h),o',s'))] + \beta \mathbb{E}[g(h \cdot (\pi(h),o',s'))] \\
  &= \alpha (T_\pi f)(h) + \beta (T_\pi g)(h).
\end{align*}
Thus \(T_\pi(\alpha f + \beta g) = \alpha T_\pi f + \beta T_\pi g\). [file:1]
\end{proof}

\begin{lemma}[Sup-norm contraction bound]
For any fixed policy \(\pi\), the operator \(T_\pi\) is a sup-norm contraction with operator norm \(1\):
\[
  \|T_\pi f\|_\infty \le \|f\|_\infty \quad \text{for all } f \in \mathcal{F}.
\] [file:1]
\end{lemma}

\begin{proof}
Let \(f \in \mathcal{F}\) and \(h \in H\). Then
\[
  |(T_\pi f)(h)|
  = \left|\mathbb{E}_{(s',o')} f(h \cdot (\pi(h),o',s'))\right|
  \le \mathbb{E}_{(s',o')} |f(h \cdot (\pi(h),o',s'))|
  \le \|f\|_\infty.
\] [file:1]
Taking the supremum over \(h\) yields \(\|T_\pi f\|_\infty \le \|f\|_\infty\). Hence \(\|T_\pi\|_{\infty \to \infty} \le 1\). [file:1]
\end{proof}

\subsection{Value Function, Discounting, and Fixed Points}

We now consider a discounted evaluation functional. Let \(r : H \to \mathbb{R}\) be an immediate reward functional and \(\gamma \in [0,1)\) a discount factor. [file:1] Define the \emph{Bellman operator}
\[
  \mathcal{B}_\pi : \mathcal{F} \to \mathcal{F}, \qquad
  \mathcal{B}_\pi f := r + \gamma T_\pi f.
\] [file:1]

\begin{lemma}[Contraction of Bellman operator]
For any fixed policy \(\pi\) and \(\gamma \in [0,1)\), the Bellman operator \(\mathcal{B}_\pi\) is a contraction in sup-norm with constant \(\gamma\):
\[
  \|\mathcal{B}_\pi f - \mathcal{B}_\pi g\|_\infty \le \gamma \|f - g\|_\infty
  \quad \text{for all } f,g \in \mathcal{F}.
\] [file:1]
\end{lemma}

\begin{proof}
For all \(h \in H\),
\begin{align*}
  \bigl|\mathcal{B}_\pi f(h) - \mathcal{B}_\pi g(h)\bigr|
  &= \bigl|r(h) + \gamma (T_\pi f)(h) - r(h) - \gamma (T_\pi g)(h)\bigr| \\
  &= \gamma \bigl|(T_\pi f)(h) - (T_\pi g)(h)\bigr| \\
  &\le \gamma \|T_\pi(f-g)\|_\infty
  \le \gamma \|f-g\|_\infty,
\end{align*}
using the previous lemma. Taking supremum over \(h\) yields the claim. [file:1]
\end{proof}

\begin{theorem}[Existence and uniqueness of value function]
For each policy \(\pi\) and \(\gamma \in [0,1)\), the Bellman equation
\[
  V_\pi = \mathcal{B}_\pi V_\pi
\]
has a unique solution \(V_\pi \in \mathcal{F}\), and value iteration
\[
  V^{(k+1)} := \mathcal{B}_\pi V^{(k)}
\]
converges in sup-norm to \(V_\pi\) for any initial \(V^{(0)} \in \mathcal{F}\). [file:1]
\end{theorem}

\begin{proof}
By the Banach fixed point theorem applied to the contraction \(\mathcal{B}_\pi\) on the complete metric space \((\mathcal{F},\|\cdot\|_\infty)\). [file:1]
\end{proof}

\subsection{Inclusion of Harm Constraints}

Let \(c : H \to \mathbb{R}_{\ge 0}\) be a per-step harm or cost functional. [file:1] Define a cumulative constrained value
\[
  J_\pi(h_0) := \mathbb{E}\left[\sum_{t=0}^{T-1} r(h_t) \,\middle|\, \pi, h_0\right], \quad
  C_\pi(h_0) := \mathbb{E}\left[\sum_{t=0}^{T-1} c(h_t) \,\middle|\, \pi, h_0\right],
\]
and consider the feasible set of policies
\[
  \Pi(\tau) := \{\pi : C_\pi(h_0) \le \tau\}.
\] [file:1]

Note that for any bounded \(r,c\), we have uniform bounds
\[
  |J_\pi(h_0)| \le T \|r\|_\infty, \qquad
  C_\pi(h_0) \le T \|c\|_\infty
\]
for all \(\pi\), hence the feasible set is nonempty for \(\tau \ge 0\). [file:1]

\begin{lemma}[Lipschitz dependence on reward and cost]
Fix \(\pi\) and horizon \(T\). Let \(r_1, r_2\) and \(c_1, c_2\) be two pairs of bounded functionals. Then
\[
  |J_{\pi,r_1}(h_0) - J_{\pi,r_2}(h_0)| \le T \|r_1 - r_2\|_\infty,
\]
\[
  |C_{\pi,c_1}(h_0) - C_{\pi,c_2}(h_0)| \le T \|c_1 - c_2\|_\infty.
\] [file:1]
\end{lemma}

\begin{proof}
By linearity of expectation and the triangle inequality:
\[
  |J_{\pi,r_1}(h_0) - J_{\pi,r_2}(h_0)|
  = \left|\mathbb{E}\sum_{t=0}^{T-1} [r_1(h_t) - r_2(h_t)]\right|
  \le \sum_{t=0}^{T-1} \|r_1 - r_2\|_\infty
  = T \|r_1 - r_2\|_\infty.
\] [file:1]
The same argument applies to \(c\). [file:1]
\end{proof}

\section{Operators on History–Belief Pairs}

\subsection{Belief State Space and Metrics}

Let \(\mathcal{B} := \Delta(S)\) be the set of probability measures on \(S\). [file:1] For \(b_1, b_2 \in \mathcal{B}\), define the total variation distance
\[
  \|b_1 - b_2\|_{\text{TV}} := \frac{1}{2} \sum_{s \in S} |b_1(s) - b_2(s)|.
\] [file:1]

Define the extended space
\[
  \mathcal{X} := H \times \mathcal{B},
\]
and consider bounded functionals \(F : \mathcal{X} \to \mathbb{R}\) with sup-norm \(\|F\|_\infty := \sup_{(h,b)\in\mathcal{X}} |F(h,b)|\). [file:1]

\subsection{Observation and Action Operators}

The draft conceptualizes operators \(\mathsf{Obs}\) and \(\mathsf{Do}\) acting on triples \((h,b,\pi)\). [file:1] For analysis, fix a policy \(\pi\) and define induced operators on functionals \(F\).

\paragraph{Observation operator.} For each observation \(o \in O\), define
\[
  (\mathsf{O}_o F)(h,b) := F(h', b'),
\]
where \(h'\) is the history extended by the observation \(o\) and \(b' := \mathsf{BayesUpdate}(b; o)\) is the Bayes-updated belief. [file:1]

\paragraph{Action operator.} For each action \(a \in A\), define
\[
  (\mathsf{A}_a F)(h,b) := \mathbb{E}_{(s',o') \sim B(\cdot \mid h,a)} F(h', b'),
\]
where \(h'\) is the history extended by \((a,o',s')\) and \(b'\) is the updated belief given \((a,o')\). [file:1]

\begin{lemma}[Norm bounds for \(\mathsf{O}_o\) and \(\mathsf{A}_a\)]
For any bounded \(F : \mathcal{X} \to \mathbb{R}\), we have
\[
  \|\mathsf{O}_o F\|_\infty \le \|F\|_\infty \quad \text{and} \quad
  \|\mathsf{A}_a F\|_\infty \le \|F\|_\infty.
\] [file:1]
\end{lemma}

\begin{proof}
For \(\mathsf{O}_o\), we have
\[
  |(\mathsf{O}_o F)(h,b)| = |F(h',b')| \le \|F\|_\infty,
\]
hence \(\|\mathsf{O}_o F\|_\infty \le \|F\|_\infty\). [file:1]

For \(\mathsf{A}_a\),
\[
  |(\mathsf{A}_a F)(h,b)|
  = \left|\mathbb{E}_{(s',o')} F(h', b') \right|
  \le \mathbb{E}_{(s',o')} |F(h',b')|
  \le \|F\|_\infty,
\]
so \(\|\mathsf{A}_a F\|_\infty \le \|F\|_\infty\). [file:1]
\end{proof}

\subsection{Lipschitz Dependence on Beliefs}

Assume the Bayes update map is Lipschitz in total variation distance: there exists \(L_{\text{Bayes}}\ge 0\) such that for any observation \(o\) and beliefs \(b_1,b_2\),
\[
  \|\mathsf{BayesUpdate}(b_1; o) - \mathsf{BayesUpdate}(b_2; o)\|_{\text{TV}} 
  \le L_{\text{Bayes}} \|b_1 - b_2\|_{\text{TV}}.
\] [file:1]

Assume also that \(F\) is Lipschitz in belief with constant \(L_F\), i.e.
\[
  |F(h, b_1) - F(h, b_2)|
  \le L_F \, \|b_1 - b_2\|_{\text{TV}}
  \quad \forall h, b_1, b_2.
\] [file:1]

\begin{lemma}[Lipschitz propagation under observation]
Under the above assumptions, \(\mathsf{O}_o F\) is Lipschitz in belief with constant \(L_F L_{\text{Bayes}}\):
\[
  |(\mathsf{O}_o F)(h, b_1) - (\mathsf{O}_o F)(h, b_2)|
  \le L_F L_{\text{Bayes}} \, \|b_1 - b_2\|_{\text{TV}}.
\] [file:1]
\end{lemma}

\begin{proof}
We have
\[
  (\mathsf{O}_o F)(h, b_i) = F(h', b'_i), \quad b'_i = \mathsf{BayesUpdate}(b_i; o), \; i=1,2.
\] [file:1]
Then
\begin{align*}
  |(\mathsf{O}_o F)(h, b_1) - (\mathsf{O}_o F)(h, b_2)|
  &= |F(h', b'_1) - F(h', b'_2)| \\
  &\le L_F \|b'_1 - b'_2\|_{\text{TV}} \\
  &\le L_F L_{\text{Bayes}} \|b_1 - b_2\|_{\text{TV}}.
\end{align*} [file:1]
\end{proof}

Analogous Lipschitz bounds can be derived for \(\mathsf{A}_a\) under appropriate smoothness of \(B\) with respect to \(b\). [file:1]

\section{Contract and Commit Operator Bounds}

\subsection{Contract Functional}

A contract is given by \(\mathcal{C} = (U, \{C_k\}_{k\in K}, B, \mathcal{P})\) as in the main text. [file:1] For each policy \(\pi \in \mathcal{P}\), define the contract value and constraint vector at initial history \(h_0\):
\[
  J_{\mathcal{C}}(\pi) := \mathbb{E}_{h \sim \mathbb{P}^{\pi,B}_{h_0}}[U(h)], \quad
  \mathbf{C}_{\mathcal{C}}(\pi) := \left( \mathbb{E}[C_k(h)] \right)_{k\in K}.
\] [file:1]

Assume:

\begin{itemize}[noitemsep]
  \item \(U, C_k\) are bounded: \(\|U\|_\infty \le M_U\), \(\|C_k\|_\infty \le M_{C_k}\). [file:1]
  \item Horizon \(T\) is finite. [file:1]
\end{itemize}

Then the same Lipschitz bounds as in the previous section imply:

\[
  |J_{\mathcal{C}}(\pi)| \le T M_U, \quad
  \|\mathbf{C}_{\mathcal{C}}(\pi)\|_\infty \le T \max_k M_{C_k}.
\] [file:1]

\subsection{Commit Operator as Argmax}

The commit operator selects a policy maximizing the contract value subject to constraints:
\[
  \mathsf{Commit}(\mathcal{C}, h_0) \in \arg\max_{\pi \in \mathcal{P}} J_{\mathcal{C}}(\pi)
  \quad \text{s.t.} \quad
  \mathbf{C}_{\mathcal{C}}(\pi) \le \boldsymbol{\tau}.
\] [file:1]

We can state continuity of the argmax correspondence under perturbations of \(U\) and \(C_k\), in a finite policy class.

\begin{theorem}[Robustness of Commit over finite policy classes]
Assume \(\mathcal{P} = \{\pi_1,\dots,\pi_N\}\) is finite. [file:1] Let \(\mathcal{C}\) and \(\tilde{\mathcal{C}}\) be two contracts that share \(B\) and \(\mathcal{P}\) but differ in \((U,\{C_k\})\) vs \((\tilde{U},\{\tilde{C}_k\})\). Define 
\[
  \delta_U := \|U - \tilde{U}\|_\infty, \quad
  \delta_{C} := \max_k \|C_k - \tilde{C}_k\|_\infty.
\] [file:1]

Then for any \(\pi \in \mathcal{P}\),
\[
  |J_{\mathcal{C}}(\pi) - J_{\tilde{\mathcal{C}}}(\pi)| \le T \delta_U, \qquad
  \|\mathbf{C}_{\mathcal{C}}(\pi) - \mathbf{C}_{\tilde{\mathcal{C}}}(\pi)\|_\infty \le T \delta_C.
\] [file:1]

In particular, if there is a margin \(\eta > 0\) such that the optimal \(\pi^*\) for \(\mathcal{C}\) satisfies
\[
  J_{\mathcal{C}}(\pi^*) \ge J_{\mathcal{C}}(\pi) + 2 T \delta_U \quad \forall \pi \ne \pi^*
\]
and the constraints remain strictly feasible under perturbations of size \(T \delta_C\), then \(\pi^*\) is also optimal for \(\tilde{\mathcal{C}}\). [file:1]
\end{theorem}

\begin{proof}
The first inequalities are immediate from the Lipschitz bounds over horizon \(T\). [file:1] The margin condition ensures that even after perturbing rewards by at most \(T\delta_U\), the ordering of \(J\) values is preserved. Constraint robustness follows similarly if \(\mathbf{C}_{\mathcal{C}}(\pi^*)\) lies strictly inside the feasible region by at least \(T\delta_C\) in sup-norm. [file:1]
\end{proof}

This provides an explicit operator norm-style robustness property for the commit operation over a finite policy class. [file:1]

\section{Norm Bounds in the Toy Implementation}

The toy environment in the main text has:

\begin{itemize}[noitemsep]
  \item \(S = \{0,1\}\), \(A = \{0,1\}\), \(O = S\). [file:1]
  \item Horizon \(T\) finite. [file:1]
  \item Reward at each step: \(r_t = \mathbf{1}\{s_t = 1\} - \lambda \mathbf{1}\{a_t = 1\}\). [file:1]
  \item Harm cost: \(c_t = \mathbf{1}\{a_t = 1\}\). [file:1]
\end{itemize}

\begin{lemma}[Explicit bounds in the toy model]
In the toy model with horizon \(T\),

\begin{enumerate}[noitemsep]
  \item For any policy \(\pi\) and any history \(h\), the cumulative reward satisfies
  \[
    -T \lambda \le U(h) \le T,
  \]
  hence \(\|U\|_\infty \le \max\{T, T\lambda\}\). [file:1]
  \item The cumulative harm cost satisfies
  \[
    0 \le C(h) \le T,
  \]
  hence \(\|C\|_\infty \le T\). [file:1]
  \item The value function \(V_\pi\) under discount \(\gamma\) satisfies the bound
  \[
    \|V_\pi\|_\infty \le \frac{\max\{1,\lambda\}}{1-\gamma}.
  \] [file:1]
\end{enumerate}
\end{lemma}

\begin{proof}
(1) At each step \(r_t \in \{-\lambda, 0, 1-\lambda, 1\}\), so \(|r_t| \le \max\{1,\lambda\}\) and the unsigned cumulative reward is at most \(T\max\{1,\lambda\}\). The signed bounds follow directly. [file:1]

(2) Harm cost is the count of actions equal to 1, so it lies in \([0,T]\). [file:1]

(3) In the infinite horizon discounted case, with per-step reward bounded by \(\max\{1,\lambda\}\) in absolute value, we have
\[
  \|V_\pi\|_\infty 
  = \sup_{h} \left|\mathbb{E}\left[\sum_{t=0}^{\infty} \gamma^t r_t \,\middle|\, h,\pi\right]\right|
  \le \sum_{t=0}^{\infty} \gamma^t \max\{1,\lambda\}
  = \frac{\max\{1,\lambda\}}{1-\gamma}.
\] [file:1]
\end{proof}

These concrete bounds show that in the toy implementation the operators and value functions are uniformly bounded and well-behaved, illustrating the general theory in a numerically tractable setting. [file:1]

\section{Summary of Operator Norm Bounds}

For reference, we summarize the main bounds established:

\begin{itemize}[noitemsep]
  \item History extension operator \(T_\pi\): \(\|T_\pi\|_{\infty \to \infty} \le 1\). [file:1]
  \item Bellman operator \(\mathcal{B}_\pi\): contraction with constant \(\gamma\) in \(\|\cdot\|_\infty\). [file:1]
  \item Observation and action operators on history–belief functionals: \(\|\mathsf{O}_o\|_{\infty \to \infty} \le 1\), \(\|\mathsf{A}_a\|_{\infty \to \infty} \le 1\). [file:1]
  \item Lipschitz propagation of belief uncertainty under \(\mathsf{O}_o\) (and analogously \(\mathsf{A}_a\)) given Lipschitz Bayes update. [file:1]
  \item Robustness of commit with respect to perturbations of contract functionals over finite policy classes, with deviation controlled by horizon and sup-norm differences in \(U\) and \(C_k\). [file:1]
\end{itemize}

These results make explicit the stability properties that were implicit in the original HCALC discussion, providing a mathematically concrete appendix suitable as an amendment to the defensive publication. [file:1]


\nocite{*}
\bibliographystyle{unsrt}  
\bibliography{references}  


\end{document}
